% Section 3: what exact means in this engine.
% Claims and status:
%  - Definitions, Lemma bf16gap, Theorem stock, Proposition stockinterval: proved here.
%  - Theorem stock's G_i for the stock cuBLAS kernel: an assumption (published model of
%    Hopper accumulation or a conservative bound); its measurement is pending (kernel).
%  - Contract witnesses: CPU tests on the repository (register entries).
\section{What ``exact'' means in this engine}\label{sec:semantics}
A certificate is a statement about a specified computation, so ``exact'' is incomplete until it names the computation whose output is preserved and the sense in which it is preserved. This section fixes both. It separates the contracts a decoder can promise, states the rule every certificate here uses, and derives the condition under which a certified token equals the token the stock head kernel would have produced.

\subsection{Five contracts}\label{sec:contracts-p1}
An exactness claim needs a reference and a sense of equality. The reference is an executable specification, not an idealized network: the weights, the kernels, the decoding policy, the random-number protocol and the execution shape, together with whether it is stock plain decoding, stock speculative verification or a recorded execution. Mathematical equivalence over the reals does not imply that two GPU reductions choose the same token, because a different batch shape or reduction order can move logits that are close~\cite{floatguide,torchnumeric}; a certificate must therefore enclose the value the reference computes, casts included. Table~\ref{tab:contracts} lists the contracts a change can claim against such a reference and the results of this paper that rest on each. They are ordered roughly from strongest to weakest, but none implies the next in general.

\begin{table}[tbp]
\small\centering
\begin{tabularx}{\linewidth}{@{}P{2.3cm}P{5.0cm}YP{1.95cm}@{}}\toprule
Contract & What must be equal & Results here that claim it & Status\\\midrule
Identical internal state & Named logical tensors (recurrent state, attention cache) at named boundaries, bitwise; not layouts or scratch buffers & Replay of recorded GDN updates (Proposition~\ref{prop:gdn}); rounding-preserving replay (P4, Section~\ref{sec:stack-state-need}); SGLang's speculative state path (H5, Section~\ref{sec:exactness-state}) & Proved; P4 and H5 tests \pendtext\\
Identical greedy outputs & Decisions, ties and stopping, with a valid state from which to continue & Certified head under R-stock (Theorem~\ref{thm:stock}, Proposition~\ref{prop:stockinterval}) & Proved under an error model; measured on replayed rows\\
Identical seeded samples & The reference's mapping from seed to output; with an inverse-CDF reference, certify with CDF bounds against the same uniform and fall back without redrawing & Certified race with the stock seeded field; fixed-noise verification against the seeded sampler on the verify pass's logits (Section~\ref{sec:fixednoise}) & Proved; measurement \pendtext\\
Exact sampling distribution & A named law: the ideal model's, the categorical law of the reference's scores, or the finite random-number implementation's; speculative correction targets the verifier's own $p$ & Races (Theorem~\ref{thm:race}); fixed-noise verification in law; residual races (Appendix~\ref{app:guards}) & Proved\\
Empirical quality equivalence & Task scores and distances within declared margins, with confidence intervals & The lossy stack under its quality budget (Section~\ref{sec:stack-stacks}) & Budget declared; measurement \pendtext\\\bottomrule
\end{tabularx}
\caption{Exactness contracts and the results that claim them. ``Measured on replayed rows'' means captured engine states replayed offline (Section~\ref{sec:rstock-share}), not a served run.}
\label{tab:contracts}
\end{table}

Small exact witnesses separate the contracts (\code{tests/test\_contracts.py}~\ev{contracts}). Reassociating an FP32 sum flips a greedy winner: $1+2^{-24}+2^{-24}$ summed left to right returns $1$ and ties a competing logit of $1$, so the first-index rule picks the competitor, while the regrouped sum returns $1+2^{-23}$ and wins, as it does in real arithmetic. Distinct FP32 values, $1$ and $1+2^{-10}$, fall in one BF16 rounding cell, so a BF16-output head ties them. Logits $(1,0)$ and $(2,0)$ have the same greedy decision but different laws: $(\tfrac23,\tfrac13)$ against $(\tfrac45,\tfrac15)$ in the base-2 mass model, and $0.731$ against $0.881$ on the first token at temperature one. Two inverse-CDF samplers that list the tokens in different orders have the same law but map the same uniform to different tokens. A state difference the current query cannot see becomes visible after one recurrent step (Section~\ref{sec:stack-state-need}), so equal outputs now do not imply equal state. Two greedy decoders that each answer two of four problems correctly have equal accuracy, yet their answers can differ on every problem, so empirical quality equivalence implies none of the stronger contracts. Equal token sequences, finally, say nothing about latency or about whether the cache that produced them is owned correctly.

For sampled verification, let $p$ be the target law at the reached prefix and $q$ the \emph{actual} proposal law, conditional on the information used to produce the proposal. A softmax buffer from a different selector is not $q$, and a deterministic proposal is a point mass. The standard one-step correction accepts $X\sim q$ with probability $\min(1,p_X/q_X)$ and otherwise samples $(p-q)_+$ normalized; the accepted mass at $i$ is $\min(p_i,q_i)$ and the residual contributes $(p_i-q_i)_+$, which sum to $p_i$~\cite{specdecode}. The certificates in this paper change how such a comparison is evaluated, never the law used in it.

\subsection{A decision needs a sound set, not a finished forward pass}\label{sec:commit}
One rule underlies every certificate in this paper. Partial computation yields information $I$ about the final logits: an approximate pass with error bounds, a transported summary, an approximate hidden state with a radius. From $I$ we keep a \emph{sound set} $Z(I)\subseteq\R^V$ that contains the reference implementation's logits. A decision rule $D(z,R)$ maps logits and the randomness fixed by the contract (a noise field, a uniform draw) to a token or an accept bit.

\begin{definition}[Commit rule]\label{def:commit}
Commit the decision $d$ when $\{D(z,R):z\in Z(I)\}=\{d\}$. Otherwise acquire more information, which must shrink or replace $Z(I)$ by another sound set, and test again.
\end{definition}
If the reference logits $z^\star$ lie in $Z(I)$, a committed $d$ equals $D(z^\star,R)$. That is the whole soundness argument, and it holds whatever produced $Z(I)$. For a box $Z(I)=\prod_j[l_j,u_j]$ the greedy rule commits $c$ when $l_c>\max_{j\neq c}u_j$. A Gumbel race with a fixed field $G$ commits $c$ when $l_c/T+G_c>\max_{j\ne c}(u_j/T+G_j)$, which needs no partition function; the field must stay fixed through every refinement~\cite{astar}. When the information is about the hidden state rather than the logits, a box is loose, and a directional certificate (Proposition~\ref{prop:directional}, Appendix~\ref{app:directional}) is the right one; Section~\ref{sec:tail} uses it.

\subsection{Agreement with the stock kernel}\label{sec:headcontracts}
Let $\mathbb B$ be the BF16 numbers. The stored head is $W\in\mathbb B^{V\times D}$ and $h\in\mathbb B^D$ is the vector the engine feeds to it. The \emph{real logit} is $z_i=\sum_jw_{ij}h_j$ in exact arithmetic. A binary format with precision $p$ (significant bits) and smallest normal exponent $e_{\min}$ has unit roundoff $u=2^{-p}$ (BF16 $p=8$, FP32 $p=24$, FP64 $p=53$). Round to nearest satisfies $|\fl(x)-x|\le u|x|+\mu$ with $\mu=2^{e_{\min}}$, which also covers flush-to-zero; directed and truncating rounding satisfy the same bound with $2u$. Write $\rnb$ for rounding to nearest-even BF16 and $\ulpb(x)$ for the BF16 spacing at $|x|$. The stock head returns $L_i=\rnb(\hat y_i)$, where $\hat y_i$ is cuBLAS's FP32 accumulation of the products.

\begin{definition}[Greedy and race contracts]\label{def:contracts}
The four references for a single decision are
\begin{align*}
\text{\textbf{R-real}:}&\quad\tok=\min\{i:z_i=\max_jz_j\},\\
\text{\textbf{R-bf16}:}&\quad\tok=\min\{i:\rnb(z_i)=\max_j\rnb(z_j)\},\\
\text{\textbf{R-stock}:}&\quad\tok=\min\{i:L_i=\max_jL_j\},\\
\text{\textbf{S-race}:}&\quad\tok=\min\{i:z_i/T+g_i=\max_j(z_j/T+g_j)\}\quad\text{for $T>0$ and a noise vector $g\in\R^V$.}
\end{align*}
\end{definition}
R-real and R-bf16 are functions of $(W,h)$ alone, so they do not depend on reduction order or batch shape. R-stock is what an unmodified server returns, and it can depend on both.

\begin{lemma}[BF16 separation]\label{lem:bf16gap}
If $y_a-y_b>\ulpb(\max(|y_a|,|y_b|))$ then $\rnb(y_a)>\rnb(y_b)$.
\end{lemma}
\begin{proof}
Rounding is nondecreasing. If both round to $v$, both lie in the preimage of $v$. When $|v|$ is not a power of two the preimage has length $\ulpb(v)$ and both values lie in the binade of $v$. When $|v|=2^e$ the preimage is $[v-\ulpb(v)/4,v+\ulpb(v)/2]$ (for $v>0$); two values on one side are within half the spacing of their binade, and values on both sides are within $\tfrac34\ulpb(v)\le\ulpb(y_a)$. At $v=0$ the preimage has the subnormal spacing. Each case contradicts the hypothesis.
\end{proof}

\begin{theorem}[Agreement with the stock head]\label{thm:stock}
Let $|\hat y_i-z_i|\le G_i$ for every $i$ and let $a$ be the R-real token. If for all $b\ne a$
\[z_a-z_b>\delta_{ab}:=G_a+G_b+\ulpb\bigl(\max(|z_a|,|z_b|)+\max(G_a,G_b)\bigr),\]
then R-stock returns $a$. If $\rnb(\hat y_i)=\rnb(z_i)$ for all $i$, R-stock equals R-bf16.
\end{theorem}
\begin{proof}
$\hat y_a-\hat y_b>\ulpb(\max(|z_a|,|z_b|)+\max(G_a,G_b))\ge\ulpb(\max(|\hat y_a|,|\hat y_b|))$, so Lemma~\ref{lem:bf16gap} gives $L_a>L_b$ for every $b$. The second claim is immediate.
\end{proof}

\paragraph{The contract used in this paper} A certified head claims R-stock at the batch shape it serves: for each row it returns the token the stock head kernel would return for the same hidden state in a batch of the same shape. It certifies a row only when the stock token is determined for every accumulator value the error model admits, by the gap condition of Theorem~\ref{thm:stock} or by the sharper interval form of Section~\ref{sec:kernel-fallback}, with $G_i$ bounding the stock kernel's pre-rounding error at that shape. For every other row it runs the stock kernel on the batch and takes that kernel's own argmax and tie rule, so those rows are the stock output by construction. Near $|z|\in[16,32)$ the gap condition asks for a margin of about $0.125+2G$.

The procedure has to establish that margin for every competitor, including rows it never re-scores. With enclosures $lo_i\le z_i\le hi_i$, define for a pair of rows
\[\hat\delta_{aj}=G_a+G_j+\ulpb\bigl(\max(|lo_a|,|hi_a|,|lo_j|,|hi_j|)+\max(G_a,G_j)\bigr).\]
A candidate $a$ is certified when every competitor $j\ne a$ satisfies $hi_j<lo_a-\hat\delta_{aj}$. Then $z_a-z_j>\hat\delta_{aj}\ge\delta_{aj}$, because $\ulpb$ is nondecreasing in the magnitude and $|z|$ is bounded by the endpoints, and Theorem~\ref{thm:stock} gives the stock token under the $G_i$ model. A competitor may therefore be excluded only when $hi_j<lo_a-\hat\delta_{aj}$, not merely when $hi_j<lo_a$. At level 0 this is checked in bulk: with $G_{\max}=\max_iG_i$, $M=\max_i\max(|lo_i|,|hi_i|)$ and $\delta_{\max}=2G_{\max}+\ulpb(M+G_{\max})$, every row with $hi_j<\max_ilo_i-\delta_{\max}$ meets the condition for whichever candidate is finally certified, since that candidate's re-scored lower bound is at least $\max_ilo_i$ (enclosures are intersected across levels) and $\hat\delta_{aj}\le\delta_{\max}$. The remaining rows are re-scored and checked pair by pair (Figure~\ref{fig:mechanism}a).

\paragraph{The assumption behind $G_i$} The condition rests on $G_i$, and $G_i$ rests on a model of the stock kernel. For an FP32 accumulation, Lemma~\ref{lem:tree} supplies $G_i$ from the reduction tree. For the tensor-core kernels cuBLAS selects on this GPU, the tree's per-node error comes from the published model of Hopper accumulation~\cite{khattak2026}, which is not a vendor contract, and a reduced-precision split-$K$ reduction admits no documented bound at all. PyTorch permits such reductions by default, so either the stock baseline disables them or the fallback must cover every shape at which cuBLAS might use them. The exact reference contains a constructed witness in which BF16 partial sums reverse a token that the FP32-model gap condition certifies (\code{tests/test\_precision.py}~\ev{precision}). The agreement of certified tokens with R-stock is therefore conditional on the stock kernel's error being within the modelled $G_i$ at the shapes served. That assumption is tested against the kernel's outputs rather than proved: on 6,005 captured plain-decode rows no certified token differed from the token the engine returned (Section~\ref{sec:rstock-share}), and the measured accumulation error of cuBLAS at the served shapes is \pend{kernel-error}.

\paragraph{Why R-bf16 is not the contract} R-bf16 is attractive: it is batch-invariant and it uses the same lowest-index rule as the stock argmax. It is not what the engine computes. The stock kernel rounds its own FP32 accumulation, which can fall on the other side of a BF16 rounding boundary from the exactly rounded value. A two-row witness shows it: with $w_0=(16,0,0)$, $w_1=(16,\tfrac1{16},2^{-30})$ and $h=(1,1,1)$, the exact logits are $16$ and $16.0625+2^{-30}$, so R-bf16 rounds them to $16$ and $16.125$ and picks token~1; a sequential FP32 accumulation loses the $2^{-30}$, rounds $16.0625$ to even, returns $16$ for both rows, and the stock argmax picks token~0 (\code{test\_bf16\_contract\_is\_not\_stock} in \code{tests/test\_precision.py}; the gap condition refuses to certify this row). Using R-bf16 as the reference would report agreement with a computation the engine does not perform. R-real remains useful as a batch-invariant reference for replay studies, and the FP32-logit contract equals the engine run with \code{-{}-enable-fp32-lm-head}; both are named as such wherever they are used.

\subsection{The interval form and the fallback}\label{sec:kernel-fallback}
The gap condition needs a bound $|\hat y_i-z_i|\le\gamma\sum_j|w_{ij}h_j|$ on the stock kernel's accumulator, and two named models supply one. The \emph{conservative} model charges $2D$ roundings of relative size $2^{-23}$ in any order; it covers any reduction tree, split-$K$ with FP32 partials and truncating adders, and gives $\gamma=6.1\times10^{-4}$ at $D=2560$. The \emph{blocked} model uses the published model of Hopper BF16 \code{wgmma} (16-term blocks, 25 fractional bits, truncation) plus an allowance for FP32 split-$K$, and gives $\gamma=1.19\times10^{-4}$~\cite{khattak2026}; it rests on that hardware model, not on vendor documentation. Neither covers split-$K$ partials stored in BF16, which PyTorch permits by default. For every batch size the engine uses, the kernel measurement therefore records which cuBLAS kernel runs and whether its output changes when reduced-precision reduction is disabled; a shape whose output changes goes to the fallback unconditionally. The same measurement records the largest observed $|\hat y_i-z_i|/\sum_j|w_{ij}h_j|$ against both models \pend{kernel-error}.

The decision can use the BF16 rounding boundaries exactly rather than a one-spacing margin. After re-scoring, each candidate has an interval $[\hat y_i^-,\hat y_i^+]$ for the stock accumulator, and because rounding is monotone its stock logit lies in $[L_i^-,L_i^+]=[\rnb(\hat y_i^-),\rnb(\hat y_i^+)]$ (Figure~\ref{fig:mechanism}b).

\begin{proposition}[Interval form of the stock decision]\label{prop:stockinterval}
Let $C$ be a set of rows such that every row $j\notin C$ has $L_j^+<\max_{i\in C}L_i^-$, and let $k=\min\{i\in C:L_i^-=\max_{j\in C}L_j^-\}$. If $L_j^+<L_k^-$ for every $j\in C$ with $j<k$ and $L_j^+\le L_k^-$ for every $j\in C$ with $j>k$, then R-stock returns $k$.
\end{proposition}
\begin{proof}
For $j<k$, $L_j\le L_j^+<L_k^-\le L_k$, so no earlier index attains the maximum. For $j>k$, $L_j\le L_j^+\le L_k^-\le L_k$, so a later index can at most tie, and the first-index rule then returns $k$. Rows outside $C$ have $L_j\le L_j^+<L_k^-\le L_k$ by the choice of $C$.
\end{proof}
Two tokens whose logits round to the same BF16 value under every admissible accumulation are therefore certified, with the tie broken by index exactly as \code{torch.argmax} breaks it; only a tie that the accumulation error could make or unmake is left undecided. This matters in practice because the gap condition has a floor: it demands a margin above one BF16 spacing even when the accumulation error is negligible, which 1.95\% of the captured positions lack (Section~\ref{sec:rstock-share}).

\paragraph{How undecided rows run} The assumption-free fallback reruns the same \code{torch.matmul} call on the whole batch at the same $M$, so cuBLAS selects the same kernel, and takes its argmax. It needs no error model, and a fallback costs time, never agreement. Its price is the whole stock head, paid by every batch that contains an undecided row, so its cost grows with the batch (Section~\ref{sec:rstock-share}). A cheaper row-wise mode relies on a property of the stock kernel that has to be measured on this cuBLAS build: that a row's BF16 logits are bitwise the same whether it is computed alone, in any subset of the batch, or against any subset of head rows. A row undecided only because of a near tie among a complete candidate list is then completed by one stock \code{matmul} over its gathered candidate rows, whose own values and tie rule decide it. That mode's claim is conditional on the measured invariance \pend{kernel-invariance}. Inside a CUDA graph each fallback is a conditional node whose predicate a routing kernel sets on the device, so a captured step has a fixed topology and pays for the fallback only when it runs (Figure~\ref{fig:mechanism}c). Section~\ref{sec:rstock-share} gives the share of undecided rows and batches for the offline replay; the same share through the kernel implementation, by error model and fallback mode, is \pend{kernel-replay}.

\subsection{Comparisons are made at the same batch shape}\label{sec:same-shape}
Stock configurations that ought to be equivalent already disagree with each other: plain decoding at different batch sizes, plain decoding against MTP speculation, and speculation at different batch sizes can emit different greedy tokens (Section~\ref{sec:divergence}). A certified head is therefore compared with the stock head at the same batch shape and in the same serving configuration, where stock is bitwise reproducible, and never with a stock run at another shape or an earlier run of another configuration. Otherwise the comparison measures the stock noise floor rather than the certificate.

\subsection{What each consumer needs}\label{sec:consumers}
Table~\ref{tab:consumers} lists the information each consumer of the head requires. A fast path is admitted by these contracts, not by a request being labelled ``decode''. Grammar masks, repetition penalties and temperature transforms are part of the scores; they cannot be applied after candidates have been discarded.

\begin{table}[tbp]
\small\centering
\begin{tabularx}{\linewidth}{@{}P{3.2cm}YP{4.1cm}@{}}\toprule
Consumer & Required information & Certificate here\\\midrule
Greedy token & Winning index under the reference kernel and tie rule & Ladder plus stock contract (Sections~\ref{sec:headcontracts}, \ref{sec:ladder})\\
Exact top-$k$ & Winning indices and scores; no partition & Ladder on membership\\
Seeded or fixed-noise sample & Argmax of perturbed scores under the keyed field & Certified race (Section~\ref{sec:race})\\
Acceptance bit & $\exp(z_X)$ against $Uq_XZ$; an interval can suffice & Guards (Appendix~\ref{app:guards}; analysis). The engine path uses fixed-noise verification instead (Section~\ref{sec:fixednoise})\\
Residual sample & The $(p-q)_+$ law, including omitted mass & Residual races (Appendix~\ref{app:guards})\\
Bonus sample & A sample from the target law under the RNG contract & Certified race\\
Top-$p$ sample & Mass above each token relative to $Z$ & Partition bounds, or fall back\\
Arbitrary log-probabilities & Requested logits and the full partition & None; ordinary path\\\bottomrule
\end{tabularx}
\caption{Consumers of the head and what each needs. The last rows are where a narrow fast path usually loses.}
\label{tab:consumers}
\end{table}
